% Open questions / future work — short. List what the catalog
% cannot answer today and where the next bits of attention should go.

We close with what the catalog cannot currently answer and where the
next bits of attention should go.

\paragraph{Fusion-kernel extraction.} The 2026-07 gap-closing
campaign against FMM-Lille (Section~\ref{sec:tables}) ended with a
sharp diagnosis: every remaining verified deficit --- six shapes, twelve
ranks in total --- traces to a single device our composition operators
do not have, the product-level \emph{fusion kernel} of
Section~\ref{sec:nonoverlap}: a handful of mixed products serving two
or more recursion slots at once. Extracting a reusable identity from
the readable specimens (one two-slot unit computes two
\(\nmpshape{2}{3}{3}\) slot-problems with 29 products instead of 30),
deciding whether the unit is self-contained or relies on whole-kernel
cancellation, and turning it into a constructor the allocation search
can elect, is the concrete next research program. Notably, disjoint-sum
($\tau$-theorem) search is \emph{not} the missing device: gaps we
initially attributed to $\tau$ (including
\(\nmpfield{\Rationals}{17}{17}{17}\), where our recombination search
reaches \(2930\)) closed via ordinary recombination, and the campaign
verified exhaustively that no remaining row is closable by existing
engines or their parameterisations. A constructive disjoint-sum
operator retains a narrower motivation: \emph{reproducing} the imported
large-format Schwartz--Zwecher schemes rather than carrying them as
opaque matrices.

\paragraph{Salt-tolerant prediction.} Any post-construction sharing
discovery would break the search-rank-equals-materialised-rank
invariant (Section~\ref{sec:nonoverlap}); if added, it belongs in a
separate polish phase after the closure converges, with the interaction
with lazy materialisation worked out first.

\paragraph{Border-rank track.} The catalog is exact-rank only. A
border-rank track would admit results like Bini's
\(\nmpfield{\Reals}{2}{2}{2}\) border-rank-\(7\) constructions and
Schönhage $\tau$-theorem upper bounds; a separate border-rank listing
already feeds the catalog browser but not yet the comparison tables.

\paragraph{Lower-bound integration.} The lower-bound registry ---
Strassen's \(\nmpfield{\Reals}{2}{2}{2} \ge 7\), Bl\"aser's family, the
recent Wang 2026 \(\nmpfield{\Ftwo}{3}{3}{3} \ge 20\) --- is not yet
surfaced in the comparison tables; a gap-to-lower-bound column (``49,
gap 4 to lower-bound 45'') would make the SOTA frontier read more
cleanly.

\paragraph{Border between catalog and hand-crafted-discovery.} When a
hand-crafted construction lands at a shape where composition was
\emph{also} predicting the same rank, the lineage should record both
discoveries and the comparison table should report both; today the
catalog picks one. A two-attribution column would force the
distinction into the table.

\paragraph{Reproducibility tooling.} The paper now prints a
date-stamped snapshot of the comparison; the live catalog browser
remains the source of truth, regenerated from the catalog and the
cited- and derived-bound registries per the conventions of Section
\ref{ssec:sota}. The remaining work item is finalising the SOTA
conventions document, so the comparison has a deterministic rule for
every corner case --- the prerequisite for the exhaustive SOTA review
that the curated paper examples wait on.
